% !TeX root = 20260818.tex
% !LW recipe=lablog XeLaTeX
\input{_support/style.tex}

\begin{document}

\reporttitle{2026-08-18}{실험 결과 및 설계 기록}

\section*{목적과 이 보고서의 범위}

이 날짜의 실험 설정, 관찰 결과, 원인에 대한 해석과 후속 판단을 함께 정리한다. 다른 날짜의 평가 방식이나 표본 수가 다른 점수를 직접 합치지 않으며, 실행 완료와 논문 근거 채택을 구별한다.

\recordnote{수합·문서 개편: 2026-10-20. 실행 날짜는 각 절과 원시 실행 목록을 따른다. 아래의 과거 계획·진행 상태는 당시 기록이며, 현재 채택 상태와 최종 결과는 후속 실험을 함께 확인한다.}

\section{2026-08-18 · 0818 — exp1 재실행 (test cap 제거), 보고서 갱신}
\label{report:20260818_0818}
\recordnote{이 절은 당시 설정의 결과다. 평가 방식 변경·정정·후속 비교는 뒤에 명시한 실험과 구별한다.}
\subsection{무엇을 바꿨나}
\texttt{test\_\allowbreak{}\allowbreak{}cap\allowbreak{}\_\allowbreak{}\allowbreak{}per\_\allowbreak{}\allowbreak{}cla\allowbreak{}ss}를 100000  \ensuremath{\rightarrow}  0으로 바꿔 exp1(100만 행 XGB vs TabPFN)을 4개 데이터셋에 재실행했다. 한 knob만 바뀐 run 쌍이 남아 있다:

\par\noindent \textbullet\enspace capped: \texttt{20260817\allowbreak{}\_\allowbreak{}\allowbreak{}072033/\allowbreak{}\allowbreak{}081129/\allowbreak{}\allowbreak{}0\allowbreak{}90418/\allowbreak{}\allowbreak{}09\allowbreak{}3806\_\allowbreak{}\allowbreak{}*\_\allowbreak{}\allowbreak{}e\allowbreak{}xp1\_\allowbreak{}\allowbreak{}1m}


\par\noindent \textbullet\enspace uncapped: \texttt{20260818\allowbreak{}\_\allowbreak{}\allowbreak{}030904/\allowbreak{}\allowbreak{}042459/\allowbreak{}\allowbreak{}0\allowbreak{}53139/\allowbreak{}\allowbreak{}07\allowbreak{}1152\_\allowbreak{}\allowbreak{}*\_\allowbreak{}\allowbreak{}e\allowbreak{}xp1\_\allowbreak{}\allowbreak{}1m}


exp2(full-XGB)도 같은 짝이 있다 — \texttt{20260817\allowbreak{}\_\allowbreak{}\allowbreak{}161427\_\allowbreak{}\allowbreak{}nfv3\_\allowbreak{}\allowbreak{}cic\allowbreak{}2018\_\allowbreak{}\allowbreak{}exp\allowbreak{}2\_\allowbreak{}\allowbreak{}fullxg\allowbreak{}b}(cap 100k)와 \texttt{20260818\allowbreak{}\_\allowbreak{}\allowbreak{}111317\_\allowbreak{}\allowbreak{}nfv3\_\allowbreak{}\allowbreak{}cic\allowbreak{}2018\_\allowbreak{}\allowbreak{}exp\allowbreak{}2\_\allowbreak{}\allowbreak{}fullxg\allowbreak{}b}(cap 없음)은 \texttt{test\_\allowbreak{}\allowbreak{}cap\allowbreak{}\_\allowbreak{}\allowbreak{}per\_\allowbreak{}\allowbreak{}cla\allowbreak{}ss} 외 args가 동일하다.

\subsection{결과}
\subsubsection{1. test cap이 소수 클래스 F1을 부풀린다 (CIC2018 full-XGB, 학습 조건 동일)}
\begingroup
\footnotesize
\begin{longtable}{>{\raggedright\arraybackslash}p{\dimexpr 0.37500\linewidth-2\tabcolsep\relax}>{\raggedleft\arraybackslash}p{\dimexpr 0.21591\linewidth-2\tabcolsep\relax}>{\raggedleft\arraybackslash}p{\dimexpr 0.20455\linewidth-2\tabcolsep\relax}>{\raggedright\arraybackslash}p{\dimexpr 0.20455\linewidth-2\tabcolsep\relax}}
\toprule
 & cap 100k & cap 없음 & Δ \\ \midrule
\endfirsthead
 & cap 100k & cap 없음 & Δ \\ \midrule
\endhead
{}\texttt{web\_\allowbreak{}\allowbreak{}atta\allowbreak{}cks} F1 (509행) & 0.603 & 0.162 & −0.441 \\
{}\texttt{infiltra\allowbreak{}tion} F1 & 0.488 & 0.429 & −0.059 \\
{}\texttt{benign} F1 & 0.876 & 0.995 & +\allowbreak{}0.119 \\
{}macro & 0.8063 & 0.7548 & −0.052 \\
{}weighted & 0.874 & 0.983 & +\allowbreak{}0.109 \\
\bottomrule
\end{longtable}
\endgroup
cap은 test 분포를 평탄화해 소수 클래스 precision을 올린다. \textbf{0812/\allowbreak{}0813 보고서의 macro 숫자는 이 회차 숫자와 직접 비교 불가.}

\subsubsection{2. exp1 (동일한 100만 행, test pool 전체) — macro F1}
\begingroup
\footnotesize
\begin{longtable}{>{\raggedright\arraybackslash}p{\dimexpr 0.20000\linewidth-2\tabcolsep\relax}>{\raggedleft\arraybackslash}p{\dimexpr 0.20000\linewidth-2\tabcolsep\relax}>{\raggedleft\arraybackslash}p{\dimexpr 0.20000\linewidth-2\tabcolsep\relax}>{\raggedleft\arraybackslash}p{\dimexpr 0.20000\linewidth-2\tabcolsep\relax}>{\raggedright\arraybackslash}p{\dimexpr 0.20000\linewidth-2\tabcolsep\relax}}
\toprule
데이터셋 & XGB & TabPFN & Δ & 시간 배수 \\ \midrule
\endfirsthead
데이터셋 & XGB & TabPFN & Δ & 시간 배수 \\ \midrule
\endhead
{}CIC2018 & 0.7568 & 0.7518 & −0.005 & 65.7× \\
{}BoT-IoT & 0.5639 & 0.5492 & −0.015 & 69.4× \\
{}ToN-IoT & 0.6964 & 0.6672 & −0.029 & 46.9× \\
{}CIC2017 & 0.6538 & 0.6505 & −0.003 & 7.9× \\
\bottomrule
\end{longtable}
\endgroup
12개 지표(4 데이터셋 × macro/\allowbreak{}weighted/\allowbreak{}tail) 중 TabPFN이 앞선 것은 2개: CIC2017 tail +\allowbreak{}0.049(전부 \texttt{heartble\allowbreak{}ed} — \textbf{test 3행}), BoT-IoT weighted +\allowbreak{}0.015 (\texttt{ddos} 0.031 \ensuremath{\rightarrow} 0.067, 양쪽 다 붕괴).

TabPFN \texttt{fit}은 데이터셋 무관 1,241–1,265초로 일정(컨텍스트가 항상 100만 행). predict만 test 행 수에 비례(56만 행 754초  \ensuremath{\rightarrow}  550만 행 7,315초).

\subsubsection{3. 100만 행  \ensuremath{\rightarrow}  전체 행 (XGB, 같은 test set)}
\begingroup
\footnotesize
\begin{longtable}{>{\raggedright\arraybackslash}p{\dimexpr 0.20000\linewidth-2\tabcolsep\relax}>{\raggedleft\arraybackslash}p{\dimexpr 0.20000\linewidth-2\tabcolsep\relax}>{\raggedleft\arraybackslash}p{\dimexpr 0.20000\linewidth-2\tabcolsep\relax}>{\raggedleft\arraybackslash}p{\dimexpr 0.20000\linewidth-2\tabcolsep\relax}>{\raggedright\arraybackslash}p{\dimexpr 0.20000\linewidth-2\tabcolsep\relax}}
\toprule
데이터셋 & 100만 & 전체 & Δ & 학습 배수 \\ \midrule
\endfirsthead
데이터셋 & 100만 & 전체 & Δ & 학습 배수 \\ \midrule
\endhead
{}CIC2018 & 0.7568 & 0.7548 & −0.002 & 12.1× \\
{}BoT-IoT & 0.5639 & 0.5610 & −0.003 & 10.2× \\
{}ToN-IoT & 0.6964 & 0.6687 & \textbf{−0.028} & 16.5× \\
{}CIC2017 & 0.6538 & 0.7643 & \textbf{+\allowbreak{}0.111} & 1.7× \\
\bottomrule
\end{longtable}
\endgroup
\par\noindent \textbullet\enspace ToN-IoT 하락은 \texttt{ransomwa\allowbreak{}re} 0.385 \ensuremath{\rightarrow} 0.150, \texttt{mitm} 0.208 \ensuremath{\rightarrow} 0.162. 795행·1,204행짜리 클래스가 물량이 늘수록 더 밀린다.


\par\noindent \textbullet\enspace CIC2017 상승 +\allowbreak{}0.111은 두 클래스가 전부다: \texttt{heartble\allowbreak{}ed}(test 3행) 0.000 \ensuremath{\rightarrow} 0.800, \texttt{infiltra\allowbreak{}tion}(test 8행) 0.000 \ensuremath{\rightarrow} 0.857. (0.800+\allowbreak{}0.857)/\allowbreak{}15 =\allowbreak{} 0.110. 100만 stratified 샘플링이 이 둘의 학습 행을 4개·12개로 깎은 것이 원인 — \textbf{데이터 양의 효과가 아니라 샘플링이 마이크로 클래스를 지운 효과.}


\subsection{결론}
\par\noindent \textbullet\enspace \textbf{TabPFN 100만 상한은 이 데이터에서 병목이 아니다.} 병목은 chronological drift와 마이크로 클래스이며, 둘 다 데이터를 더 넣어서 해결되지 않는다.


\par\noindent \textbullet\enspace \textbf{동일 데이터 조건에서 TabPFN은 4/\allowbreak{}4 데이터셋 macro F1이 XGBoost보다 낮다.} 0812/\allowbreak{}0813의 "명확히 이기는 설정 없음"이 test cap 제거 후에도 유지되며, 오히려 더 명확해졌다.


\par\noindent \textbullet\enspace \textbf{stratified ratio-preserving 샘플링은 마이크로 클래스에 대해 안전하지 않다.} 클래스당 최소 행 보장이 없으면 학습 행이 한 자리수까지 떨어진다.


\subsection{다음에 할 것}
\par\noindent \textbullet\enspace [ ] \texttt{heartble\allowbreak{}ed}/\allowbreak{}\texttt{sql-\allowbreak{}inje\allowbreak{}ction}/\allowbreak{}\texttt{infiltra\allowbreak{}tion}(test 3·5·8행)을 macro F1 집계에서 분리 보고할지 결정. 현행은 15클래스 macro가 3행짜리 클래스에 1/\allowbreak{}15 =\allowbreak{} 0.067을 준다.


\par\noindent \textbullet\enspace [ ] 학습 샘플링에 클래스당 최소 행(예: min 1000 또는 전량) 옵션 추가 — 새 플래그, 새 run.


\par\noindent \textbullet\enspace [ ] 0817의 미해결 항목(train+\allowbreak{}val 컨텍스트 경로)은 그대로 남아 있다.


\subsection{Files}
\par\noindent \textbullet\enspace 보고서: \texttt{docs/\allowbreak{}\allowbreak{}tab\allowbreak{}pfn\_\allowbreak{}\allowbreak{}repo\allowbreak{}rt.\allowbreak{}html} (이 회차 반영 완료)

\par\noindent \textbullet\enspace 온라인 아티팩트: https:/\allowbreak{}/\allowbreak{}claude.ai/\allowbreak{}code/\allowbreak{}artifact/\allowbreak{}a586cfcc-0f5d-4c0d-972e-2b727d6880ce


\par\noindent \textbullet\enspace \textbf{주의}: 온라인 쪽에만 있던 "test 샘플 수 내림차순 정렬" 개정판과 로컬 쪽 exp1 추가분이 따로 놀고 있었다. 온라인을 베이스로 병합해 양쪽을 같은 내용으로 맞춰 놓았다. exp1도 실험 2와 같은 형식으로 데이터셋별 클래스별 F1 표(1.1–1.4)를 넣었다.



\par\noindent \textbullet\enspace run: \texttt{tabpfn/\allowbreak{}\allowbreak{}r\allowbreak{}esults/\allowbreak{}\allowbreak{}2\allowbreak{}0260818\_\allowbreak{}\allowbreak{}*\_\allowbreak{}\allowbreak{}exp1\_\allowbreak{}\allowbreak{}1\allowbreak{}m}, \texttt{tabpfn/\allowbreak{}\allowbreak{}r\allowbreak{}esults/\allowbreak{}\allowbreak{}2\allowbreak{}0260818\_\allowbreak{}\allowbreak{}*\_\allowbreak{}\allowbreak{}exp2\_\allowbreak{}\allowbreak{}f\allowbreak{}ullxgb}


\par\noindent \textbullet\enspace 이전 회차(cap 100k): \texttt{tabpfn/\allowbreak{}\allowbreak{}r\allowbreak{}esults/\allowbreak{}\allowbreak{}2\allowbreak{}0260817\_\allowbreak{}\allowbreak{}*}




\section*{출처와 실행 추적}

상세 수치의 원시 실행은 \path{inventory/runs.csv}, 실험별 연결은 \path{inventory/experiments.csv}에서 확인한다. 표·그림의 재현 입력은 \path{evidence/}와 \path{data/}에 둔다. 문서 이동·통합 내역은 \path{configs/experiments/migrations/20261020_documents/manifest.json}에 기록했다.

\end{document}
